% Standalone XeLaTeX source. Copy to main.tex in Overleaf; select XeLaTeX.
\documentclass[11pt,a4paper]{article}
\usepackage[margin=23mm,headheight=15pt]{geometry}
\usepackage{fontspec}
% CMU fonts are supplied by standard full TeX Live / Overleaf installations.
% The fallback supports the lean local audit runtime without extra uploads.
\IfFontExistsTF{CMU Serif}
  {\setmainfont{CMU Serif}}
  {\setmainfont{NimbusRoman-Regular.otf}[BoldFont=NimbusRoman-Bold.otf,ItalicFont=NimbusRoman-Italic.otf,BoldItalicFont=NimbusRoman-BoldItalic.otf]}
\IfFontExistsTF{CMU Sans Serif}
  {\setsansfont{CMU Sans Serif}}
  {\setsansfont{DejaVu Sans}}
\IfFontExistsTF{CMU Typewriter Text}
  {\setmonofont{CMU Typewriter Text}[Scale=0.82]}
  {\setmonofont{DejaVu Sans Mono}[Scale=0.82]}
\usepackage{amsmath,amssymb,amsthm,mathtools}
\usepackage{booktabs,longtable,array,enumitem}
\usepackage{xcolor}
\usepackage{hyperref}
\hypersetup{colorlinks=true,linkcolor=blue!50!black,urlcolor=blue!60!black,citecolor=blue!50!black,pdftitle={Решётки рангов 9 и 10: новые результаты 7 октября 2026}}
\setcounter{tocdepth}{1}\hypersetup{bookmarksdepth=2}
\usepackage{fancyhdr}
\usepackage{needspace}
\pagestyle{fancy}\fancyhf{}
\fancyhead[L]{\small Решётки рангов 9 и 10}
\fancyhead[R]{\small 7 октября 2026}
\fancyfoot[C]{\thepage}
\setlength{\parindent}{0pt}\setlength{\parskip}{5pt}
\setlength{\emergencystretch}{3em}
\raggedright
\allowdisplaybreaks
\setlist{nosep,topsep=4pt,leftmargin=*}
\newtheorem{theorem}{Теорема}[section]
\newtheorem{lemma}[theorem]{Лемма}
\newtheorem{proposition}[theorem]{Предложение}
\newtheorem{corollary}[theorem]{Следствие}
\theoremstyle{remark}\newtheorem{remark}[theorem]{Замечание}
\renewcommand{\proofname}{Доказательство}
\renewcommand{\contentsname}{Содержание}
\renewcommand{\refname}{Литература и исходные материалы}
\newcommand{\Z}{\mathbb Z}\newcommand{\R}{\mathbb R}\newcommand{\F}{\mathbb F}
\newcommand{\GL}{\operatorname{GL}}\newcommand{\rank}{\operatorname{rank}}
\newcommand{\spanR}{\operatorname{span}_{\R}}\newcommand{\diag}{\operatorname{diag}}
\newcommand{\mink}{\operatorname{min}}
\newcommand{\Kalt}{K_{10}^{\mathrm{alt}}}
\newcommand{\file}[1]{{\small\nolinkurl{#1}}}
\newcolumntype{P}[1]{>{\raggedright\arraybackslash}p{#1}}
\begin{document}
\begin{titlepage}
\vspace*{18mm}
{\Large\sffamily Конфиденциальный исследовательский отчёт\par}
\vspace{14mm}
{\huge\bfseries Чётные целые решётки\par рангов 9 и 10\par}
\vspace{9mm}
{\Large Точные минимумы, случаи равенства\par и кластерные расширения\par}
\vspace{12mm}
{\large Новые результаты вычислений и внутренней проверки\par 7 октября 2026 года\par}
\vspace{14mm}
\textbf{Основные выводы.}
\[
e_9(6)=21870,\qquad \text{минимизатор единственен с точностью до изометрии};
\]
\[
e_{10}(6)=54675.
\]
В подклассе ранга 10, содержащем $\sqrt3E_8$, минимум равен $54675$ и достигается ровно двумя неизометричными решётками. Для обеих все нетривиальные антиподные кластеры с квадратом диаметра $\beta<6$ ухудшают плотность при гарантированном радиусе. В ранге 9 отсутствие кластерного улучшения установлено для всех сечений, содержащих $\sqrt3E_8$, без ограничения на их определитель.

\vfill
\textbf{Статус текста.} Это дополнение к полному отчёту аудита от 10 сентября 2026 года. Точные конечные сертификаты выполнены; их логические зависимости указаны явно. Глобальный минимум ранга 10 установлен. Полная классификация его равенства и всех вложений в ранг 48 остаётся открытой. Внешнего рецензирования и подтверждения научной новизны пока нет.
\end{titlepage}
{\small\tableofcontents}
\newpage

\section{Какие вопросы закрыты и в каком смысле}

Исходная работа посвящена сечениям экстремальных чётных унимодулярных решёток ранга 48 и упаковкам в размерностях 38--43. Предыдущая ветка дала компьютерно поддержанную нижнюю оценку $e_9(6)\ge21870$ и явный пример $K_9$, достигающий равенства. Настоящий отчёт фиксирует \emph{добавочные} результаты текущего исследования.

\begin{longtable}{P{0.06\linewidth}P{0.37\linewidth}P{0.49\linewidth}}
\toprule
№ & Поставленный вопрос & Установленное положение\\\midrule
\endhead
1 & Единственность решётки определителя $21870$ & Закрыт для абстрактных чётных целых решёток ранга 9 с минимумом не меньше 6: единственный класс изометрии $K_9$.\\[5pt]
2 & Точное $e_{10}(6)$ & Значение установлено: $e_{10}(6)=54675$. В подклассе с $\sqrt3E_8$ ровно два типа равенства; вне этого подкласса равенство ещё не классифицировано.\\[5pt]
3 & Примитивные вложения в заданные решётки ранга 48 & Над фиксированным напечатанным $K_9\subset P_{48p}$ имеется ровно 19 минимальных примитивных надсечений ранга 10; все типа $K_{10}$, ни одного альтернативного типа. Все орбиты вложений произвольного $K_9$ не перечислены.\\[5pt]
4 & Оптимальность кластерных расширений & Закрыта для всех $\beta<6$ над обоими минимальными типами ранга 10, содержащими $\sqrt3E_8$. В ранге 9 закрыта для любого определителя при наличии $\sqrt3E_8$. Для любых определителей ранга 10 с $\sqrt3E_8$ получены также частичные равномерные оценки приложения~\ref{app:uniform10}. Общая задача открыта.\\[5pt]
5 & Внешняя проверка и научная новизна & Выполнены отдельные внутренние проверки и обзор первичных источников. Внешний экспертный вердикт и доказательство новизны не получены.\\
\bottomrule
\end{longtable}

Термин ``закрыт'' в пунктах 1, 2 и 4 означает наличие приведённого математического аргумента вместе с точными конечными сертификатами в указанной области применимости. Это не означает формализацию в proof assistant или независимое человеческое рецензирование.

\section{Обозначения и общие леммы}

\subsection{Определитель, минимум и двойственная решётка}

Решётка $L$ ранга $r$ задаётся базисом в евклидовом пространстве; его матрица Грама $G$ положительно определена. Полагаем
\[
d(L)=\det G,\qquad m(L)=\min_{0\ne x\in\Z^r}x^{\mathsf T}Gx.
\]
Таким образом, $d(L)$ есть квадрат коволюма, а $m(L)$ --- квадрат минимальной длины. Решётка целая, если все скалярные произведения целые, и чётная, если все нормы чётные. В базисе это условия $G_{ij}\in\Z$, $G_{ii}\in2\Z$.

Двойственная решётка
\[
L^*=\{y\in\spanR L:(y,L)\subseteq\Z\}
\]
имеет определитель $d(L^*)=d(L)^{-1}$. Для целой $L$ конечная группа $A_L=L^*/L$ имеет порядок $d(L)$. Определим
\[
\mu_L(a)=\min_{y\in a}\|y\|^2,\qquad
R_L=\max_{a\in A_L}\mu_L(a).
\]
$R_L$ --- максимум минимумов \emph{дискретных} классов; это не квадрат обычного радиуса покрытия всего пространства.

Нас интересует
\[
e_r(6)=\min\{d(L):\rank L=r,\ L\text{ чётная целая},\ m(L)\ge6\}.
\]
Матрицы $G,G'$ изометричны над $\Z$, если $G'=U^{\mathsf T}GU$ при $U\in\GL_r(\Z)$.

\subsection{Примитивные сечения и флаги}

Подрешётка $J\subset L$ примитивна, если $J=L\cap\spanR J$, эквивалентно отсутствию кручения в $L/J$. Для примитивной подрешётки ортогональная проекция $P$ решётки $L$ на $\spanR(J)^\perp$ является решёткой и
\begin{equation}\label{eq:projectiondet}
d(L)=d(J)d(P).
\end{equation}
Действительно, базис $J$ продолжается до базиса $L$. Определитель блочной матрицы Грама равен произведению определителя блока $J$ на определитель его дополнения Шура; последний есть матрица Грама проекций добавленных векторов.

Если $y$ --- кратчайший ненулевой вектор $L^*$, то $y$ примитивен и
\begin{equation}\label{eq:dualflag}
J=L\cap y^\perp,\qquad d(J)=d(L)\|y\|^2=d(L)m(L^*).
\end{equation}
Примитивность $y$ означает $(y,L)=\Z$, поэтому шаг проекции на его прямую равен $1/\|y\|$. Формула следует из \eqref{eq:projectiondet}.

Постоянная Эрмита определяется для \emph{всех вещественных евклидовых решёток}:
\[
\gamma_r=\sup_L\frac{m(L)}{d(L)^{1/r}}.
\]
По \eqref{eq:dualflag} последовательные сечения удовлетворяют
\begin{equation}\label{eq:windows}
d_{r-1}\le\gamma_r d_r^{(r-1)/r},\qquad
d_r\ge6^r/\gamma_r^r.
\end{equation}
В исходном доказательстве использованы классические значения
\[
(\gamma_r^r)_{r=1}^8=(1,4/3,2,4,8,64/3,64,256)
\]
и неравенство Морделла
\begin{equation}\label{eq:mordell}
\gamma_r\le\gamma_{r-1}^{(r-1)/(r-2)}.
\end{equation}
Для новых оценки разрыва и общей границы ранга 10 дополнительно импортируется публичная теорема $\gamma_9=2$ из \cite{dw}. Она не нужна для новой классификации случая $d_9=21870$.

Для чётных матриц нечётного порядка определитель чётен; для порядка $2m$ он сравним с $0$ или $(-1)^m$ по модулю 4. Эти ограничения учитываются при округлении всех окон.

\subsection{Полный критерий одношагового расширения}

Пусть $J$ чётная целая и $m(J)\ge6$. Примитивное расширение на один ранг имеет вид
\[
K=J+\Z(y+z),\qquad y\in J^*,\quad z\perp J,\quad s=\|z\|^2>0.
\]
Выберем $y$ кратчайшим в классе $a=y+J$. Тогда
\begin{align}
d(K)&=d(J)s,\label{eq:extensiondet}\\
K\text{ чётная целая}&\Longleftrightarrow s+\mu_J(a)\in2\Z,\label{eq:extensioneven}\\
m(K)\ge6&\Longleftrightarrow
m^2s+\mu_J(ma)\ge6\quad(m\in\Z\setminus\{0\}).\label{eq:extensionmin}
\end{align}
Последняя формула получается минимизацией по каждому слою $J+m(y+z)$. Проверять достаточно положительные $m$ с $m^2s<6$. Если класс имеет порядок $h$, слой $m=h$ также обязателен: он даёт условие $h^2s\ge6$.

Из слоя $m=1$ следует полезная оценка
\begin{equation}\label{eq:gluebound}
d(K)\ge d(J)(6-R_J).
\end{equation}

Если расширение продолжает именно флаг кратчайших двойственных нормалей, то
\begin{equation}\label{eq:lift}
s\ge\frac{3}{4m(J^*)}.
\end{equation}
В самом деле, кратчайший вектор $u\in J^*$ можно поднять в $K^*$, выбрав коэффициент вдоль перпендикулярной двойственной нормали в $[-1/2,1/2]$. Его норма не превосходит $m(J^*)+1/(4s)$, тогда как кратчайшая нормаль имеет норму $1/s$. Отсюда \eqref{eq:lift}. В цепочке рангов 6, 7, 8 это даёт
\begin{equation}\label{eq:lift678}
d_8\ge\frac{3d_7^2}{4d_6}.
\end{equation}

\section{Единственность решётки ранга 9}

\begin{theorem}\label{thm:unique}
Если $L$ --- положительно определённая чётная целая решётка ранга 9, $m(L)\ge6$ и $d(L)=21870$, то $L\cong K_9$. В соединении с прежней нижней оценкой это классифицирует все минимизаторы $e_9(6)$: имеется ровно один класс изометрии.
\end{theorem}

\subsection{Почему требовался новый анализ границы}

Прежний перебор исходил из $d_9\le21868$. Равенство $21870$ меняет верхние окна, поэтому перенести старый вывод без дополнительных проверок было бы ошибкой. Используя $\gamma_9\le2^{8/7}$, получаем
\[
\left\lfloor2^{8/7}21870^{8/9}\right\rfloor=15910.
\]
Это проверено целочисленным сравнением
\[
15910^{63}\le2^{72}21870^{56}<15911^{63}.
\]
Так как $d_8\equiv0,1\pmod4$, верхняя допустимая граница равна \textbf{15909}, а не 15908.

\begin{center}
\begin{tabular}{c*{8}{r}}
\toprule
Ранг &8&7&6&5&4&3&2&1\\\midrule
Сырой предел &15910&9494&4648&1894&634&178&39&7\\
С учётом сравнений &15909&9494&4648&1894&633&178&39&6\\
\bottomrule
\end{tabular}
\end{center}

Новые случаи возникают только при $d_6=4648$, $d_7=9494$, $d_8=15909$. Кроме того, возвращается случай $d_8=6561$, ранее исключённый для строгого контрпримера. Равенство в теореме Эрмита ранга 8 и единственность её евклидова экстремизатора дают тогда $L_8\cong\sqrt3E_8$.

\subsection{Исключение трёх новых концов окон}

Точный генератор расширений \eqref{eq:extensiondet}--\eqref{eq:extensionmin} выполнен на всех прежних представителях родителей:
\begin{center}
\begin{tabular}{lrl}
\toprule
Целевой случай & Число родителей & Результат\\\midrule
Ранг 6, определитель 4648 &118&Ни одного расширения\\
Ранг 7, определитель 9494 &138&Ни одного расширения\\
Ранг 8, определитель 15909 &6&Ни одного расширения\\
\bottomrule
\end{tabular}
\end{center}
Полнота второго списка следует из пустоты первого нового случая. Для третьего списка заново применено \eqref{eq:lift678} ко всем 2235 прежним записям ранга 7 при увеличенном пределе 15909. Остаются те же 96 записей, сведённые точными унимодулярными матрицами к шести представителям.

Эти 262 граничных исключения дополнительно проверены независимым алгоритмом, который не вызывает прежний перечислитель коротких векторов. Для 188 случаев найдены явные двойственные векторы, нарушающие \eqref{eq:lift}; 74 случая исключены конечной квадратичной формой дискриминантной группы.

Опишем второй метод. Для матрицы родителя $G$, $d=\det G$, $A=\operatorname{adj}G$ класс с колонкой скалярных произведений $p$ кодируется
\[
c=Ap\pmod d,\qquad Q(c)=p^{\mathsf T}Ap\pmod{2d}.
\]
$Q$ корректно определена, поскольку замена $p$ на $p+Gz$ меняет числитель на
\[
2d\,p^{\mathsf T}z+d\,z^{\mathsf T}Gz\in2d\Z.
\]
Все $d$ классов перечисляются сложением колонок $A$ с рекуррентной формулой
\[
Q(c+Ae_i)=Q(c)+2c_i+A_{ii}\pmod{2d}.
\]
Расширение определителя $T$ с чётной последней диагональю $n$ требует
\[
T=dn-p^{\mathsf T}Ap,\qquad Q(c)\equiv-T\pmod{2d}.
\]
Если такого значения нет, расширение невозможно вне зависимости от минимума. Для двойственного свидетельства достаточно ненулевого $p$, удовлетворяющего
\[
4T(p^{\mathsf T}Ap)<3d^2.
\]
Это строго противоречит \eqref{eq:lift}, применённой даже к данному, не обязательно кратчайшему, двойственному вектору.

\subsection{Три возможных родителя ранга 8}

Прежняя полная классификация шести ветвей и \eqref{eq:gluebound} оставляют для равенства только
\[
J_0=\sqrt3E_8,\qquad J_1=J_{14580},\qquad J_2=J_{15309a}.
\]
Для остальных родителей нижние границы равны 32076, 36504, 39930, 42282 или вообще отсутствует допустимое продолжение флага.

Чтобы зафиксировать два менее стандартных родителя, обозначим через $H$ матрицу $\sqrt3E_7$ из старого сертификата и через
\[
\operatorname{border}(H,b,n)=\begin{pmatrix}H&b\\b^{\mathsf T}&n\end{pmatrix}.
\]
Тогда в зафиксированном базисе сертификата
\begin{align*}
J_1&=\operatorname{border}(H,(2,-3,-1,-3,-3,-3,-3)^{\mathsf T},6),\\
J_2&=\operatorname{border}(H,(0,2,2,-1,-1,-1,-1)^{\mathsf T},6).
\end{align*}
Здесь буквами решёток обозначены и их матрицы; $H$ равна верхнему левому блоку размера $7\times7$ матрицы $G_8$, выписанной в разделе 4. Все замены базисов входят в машиночитаемый пакет. Новое доказательство использует прежнюю полную классификацию как явно указанный вход.

Для каждого родителя фиксировано $s=21870/d(J)$ и заново вычислена полная таблица минимумов дискриминантных классов:
\begin{center}
\begin{tabular}{lrrrr}
\toprule
Родитель &$s$&$R_J$&Число классов&Допустимые классы\\\midrule
$\sqrt3E_8$ &$10/3$&$8/3$&6561&1920\\
$J_{14580}$ &$3/2$&$9/2$&14580&1\\
$J_{15309a}$ &$10/7$&$32/7$&15309&2\\
\bottomrule
\end{tabular}
\end{center}
Полные перечисления обследовали соответственно 26641, 190491 и 227617 двойственных векторов. Полнота таблицы удостоверяется наличием $d(J)$ различных классов внутри точно перечисленного шара. Все векторы вне шара длиннее найденных представителей. Для $J_2$ слой $m=2$ действительно проверяется, поскольку $4s=40/7<6$.

\subsection{Изометрии всех 1923 презентаций}

Пусть $G_8=3C$ --- ведущий блок $K_9$, где $C$ является матрицей $E_8$. Отражения в восьми базисных корнях
\[
R_i=I-e_i(\text{строка }i\text{ матрицы }C)
\]
целочисленны, имеют определитель $-1$ и сохраняют $C$. Обход орбиты одного допустимого класса даёт ровно 1920 различных классов, то есть все допустимые классы $J_0$. Сохранено дерево этой орбиты.

Изометрия родителя переносит клеящий вектор в тот же класс, что и выбранный представитель следующего расширения. Разность представителей лежит в $J_0$; её добавление к последнему базисному вектору является унимодулярным изменением базиса. Поэтому все 1920 расширений изометричны. Независимый проверяющий алгоритм явно построил и перемножил 1919 таких блочных матриц, по одной для каждого ребра дерева. Не требуется утверждать, что выбранные отражения порождают всю группу автоморфизмов.

Три оставшихся расширения имеют колонки
\begin{align*}
J_1:&\quad(-2,0,1,0,0,0,0,-3)^{\mathsf T},\\
J_2:&\quad(3,-2,-2,-2,-2,-2,-2,-3)^{\mathsf T},\\
J_2:&\quad(0,-3,-3,0,0,0,0,-3)^{\mathsf T},
\end{align*}
и последнюю диагональ 6. Для каждой сохранена целая матрица $U$ с
\[
\det U=\pm1,\qquad U^{\mathsf T}K_9U=G.
\]
Эти тождества проверены точно двумя процедурами. Получен один класс изометрии, что завершает доказательство теоремы \ref{thm:unique} относительно исходной нижнеранговой классификации.

\section{Ранг 10 при наличии \texorpdfstring{$\sqrt3E_8$}{sqrt(3)E8}}

\subsection{Фиксированные матрицы}

Для самостоятельного использования выпишем общий блок всех трёх основных примеров:
\[
G_8=\begin{pmatrix}
6&-3&-3&-3&-3&-3&-3&3\\
-3&6&3&3&3&3&3&0\\
-3&3&6&0&3&0&0&-3\\
-3&3&0&6&3&3&3&0\\
-3&3&3&3&6&3&3&-3\\
-3&3&0&3&3&6&3&0\\
-3&3&0&3&3&3&6&0\\
3&0&-3&0&-3&0&0&6
\end{pmatrix},\qquad C=G_8/3.
\]
Положим
\[
p=(-2,1,3,-1,1,0,1,-1)^{\mathsf T},\qquad
q=(1,-1,1,0,-1,-2,-2,1)^{\mathsf T}.
\]
Тогда
\begin{equation}\label{eq:examples}
K_9=\begin{pmatrix}G_8&p\\p^{\mathsf T}&6\end{pmatrix},\quad
K_{10}=\begin{pmatrix}G_8&p&q\\p^{\mathsf T}&6&2\\q^{\mathsf T}&2&6\end{pmatrix},\quad
\Kalt=\begin{pmatrix}G_8&p&-p\\p^{\mathsf T}&6&-1\\-p^{\mathsf T}&-1&6\end{pmatrix}.
\end{equation}
Матрица $C$ чётная, положительно определённая и унимодулярная. Её ведущие главные миноры равны $2,3,4,4,4,3,2,1$. Следовательно, она задаёт $E_8$.

\subsection{Классы \texorpdfstring{$E_8/3E_8$}{E8/3E8}}

Класс $J^*/J$ для $J=\sqrt3E_8$ можно кодировать $u\in E_8/3E_8$; его минимум равен одной трети минимальной $E_8$-нормы данного остатка. Полное перечисление до нормы 8 даёт:
\begin{center}
\begin{tabular}{crrrrr}
\toprule
Минимальная $E_8$-норма &0&2&4&6&8\\\midrule
Число классов &1&240&2160&2240&1920\\
\bottomrule
\end{tabular}
\end{center}
Сумма равна $3^8=6561$. Полные слои векторов этих норм содержат $1,240,2160,6720,17520$ векторов. Поэтому
\begin{equation}\label{eq:e8radius}
R_{\sqrt3E_8}=8/3.
\end{equation}
Назовём классы минимальной $E_8$-нормы 8 глубокими; это всего лишь обозначение конечного набора из 1920 классов.

\begin{theorem}\label{thm:rank10e8}
Пусть $L$ чётная целая решётка ранга 10, $m(L)\ge6$, содержащая подрешётку $J\cong\sqrt3E_8$. Тогда $d(L)\ge54675$. Равенство достигается ровно двумя классами изометрии: $K_{10}$ и $\Kalt$ из \eqref{eq:examples}.
\end{theorem}

\begin{proof}[Доказательство нижней оценки]
Сначала $J$ автоматически примитивна. Её насыщение $\overline J=L\cap\spanR J$ лежит в $J^*$. При ненулевом классе $\overline J/J$ вычитание вектора $J$ даёт ненулевой вектор нормы не больше $8/3$, вопреки $m(L)\ge6$.

Проекция $P$ на двумерное ортогональное дополнение является решёткой. Для любого ненулевого $z\in P$ выберем подъём $y+z\in L$. Целочисленность даёт $y\in J^*$; вычитая $J$, добиваемся $\|y\|^2\le8/3$. Отсюда
\[
6\le\|y\|^2+\|z\|^2\le8/3+\|z\|^2,
\qquad m(P)\ge10/3.
\]
Для любой двумерной решётки
\[
d(P)\ge\frac34m(P)^2.
\]
Чтобы увидеть это, берём кратчайший примитивный вектор $z_1$, продолжаем до базиса и редуцируем $z_2$ так, чтобы $|(z_1,z_2)|\le\|z_1\|^2/2$. При $a=\|z_1\|^2$, $b=\|z_2\|^2\ge a$ определитель не меньше $ab-a^2/4\ge3a^2/4$. Таким образом,
\[
d(L)=6561d(P)\ge6561\cdot\frac34\left(\frac{10}{3}\right)^2=54675.
\]
\end{proof}

\subsection{Полная классификация равенства}

Равенство вынуждает, после возможной смены знака базиса, матрицу проекции
\begin{equation}\label{eq:binaryprojection}
S=\begin{pmatrix}10/3&5/3\\5/3&10/3\end{pmatrix}.
\end{equation}
Её минимальные направления --- $\pm z_1$, $\pm z_2$, $\pm(z_1-z_2)$. Пусть $a,b\in E_8/3E_8$ --- классы подъёмов. Минимум 6 требует, чтобы $a$, $b$ и $a-b$ были глубокими.

Обратно, всякая такая пара задаёт допустимую решётку. Если $u,v$ --- представители норм 8, её матрица равна
\begin{equation}\label{eq:pairmatrix}
\begin{pmatrix}
3C&Cu&Cv\\
u^{\mathsf T}C&6&(u^{\mathsf T}Cv+5)/3\\
v^{\mathsf T}C&(u^{\mathsf T}Cv+5)/3&6
\end{pmatrix}.
\end{equation}
Из нормы $u-v\equiv8\pmod6$ следует $u^{\mathsf T}Cv\equiv1\pmod3$, поэтому матрица целая и чётная. Шесть минимальных проекций дают полную норму 6; всякий другой ненулевой вектор проекции имеет норму не меньше 10. Внутри $J$ минимум равен 6. Следовательно, \eqref{eq:pairmatrix} имеет минимум 6 и определитель 54675.

Изменение представителей на $3E_8$ добавляет векторы $J$ к подъёмам; это целочисленная замена базиса. Изометрия $E_8$ также переносится на расширение. Поэтому для верхней оценки числа классов достаточно перечислить орбиты допустимых пар под любой явно построенной подгруппой автоморфизмов.

Все 240 корней $E_8$ перечислены; после отождествления противоположных корней получается 120 отражений. Их целочисленность, определитель $-1$ и сохранение $C$ проверены точно. Они транзитивны на 1920 глубоких классах. Зафиксируем
\[
a=w\pmod3,\qquad w=C^{-1}p=(-1,-1,2,-1,0,0,1,1)^{\mathsf T}.
\]
Имеется ровно 561 глубокий класс $b$ с глубоким $a-b$. Из 120 отражений 36 фиксируют $a$. Подгруппа, ими порождённая, разбивает 561 класс на орбиты
\[
560+1;
\]
одноточечная орбита --- $b=-a$, большая содержит исходный $K_{10}$. Значит, полных решёток не больше двух типов. Не требуется доказывать, что эти 36 отражений порождают весь стабилизатор.

Оба типа существуют по \eqref{eq:examples} и неизометричны:
\[
\rank_{\F_3}K_{10}=4,\qquad \rank_{\F_3}\Kalt=3.
\]
Целая изометрия индуцирует невырожденную замену базиса по модулю 3 и сохраняет ранг. Это различает абстрактные решётки независимо от выбранного вложения $E_8$ и завершает классификацию равенства.

\subsection{Дискриминантные группы и слои векторов}

Для обоих типов $d=54675=3^7\cdot5^2$ и ранг по модулю 5 равен 8. Ранги по модулю 3 и сумма $3$-адических показателей дают соответственно показатели $(1,1,1,1,1,2)$ и $(1,1,1,1,1,1,1)$. Поэтому
\begin{align*}
A_{K_{10}}&\cong(\Z/3\Z)^4\oplus\Z/15\Z\oplus\Z/45\Z,\\
A_{\Kalt}&\cong(\Z/3\Z)^5\oplus(\Z/15\Z)^2.
\end{align*}
Точные числа ориентированных векторов малых норм:
\begin{center}
\begin{tabular}{lrrr}
\toprule
Решётка &Норма 6&Норма 8&Норма 10\\\midrule
$K_{10}$&294&216&756\\
$\Kalt$&294&216&762\\
\bottomrule
\end{tabular}
\end{center}
Равенство первых двух слоёв само по себе не означало бы изометрию; различие рангов по модулю 3 даёт более короткое свидетельство неизометричности.

Дополнительное описание альтернативы: $\Kalt$ является чётной надрешёткой индекса 3 решётки $\sqrt3E_8\perp\sqrt5A_2$. Два ортогональных вектора в координатах \eqref{eq:examples}, $(0,\ldots,0,1,1)$ и $(-w,1,-2)$, имеют матрицу
\[
\begin{pmatrix}10&-5\\-5&10\end{pmatrix}.
\]
Определитель прямой суммы равен $6561\cdot75$, а деление на квадрат индекса $9$ даёт 54675. Ненулевые клеящие классы имеют минимумы $8/3$ и $10/3$, сумма которых равна 6.

\section{Разрыв в ранге 9 и предварительная граница ранга 10}

\begin{proposition}\label{prop:gap}
При использовании публичной теоремы $\gamma_9=2$ всякая чётная целая решётка ранга 9 с $m(L)\ge6$ и $d(L)\le24450$ изометрична $K_9$. В частности, для $L\not\cong K_9$ имеем $d(L)\ge24452$.
\end{proposition}

\begin{proof}
Прежняя нижняя оценка даёт $D=d(L)\ge21870$. Для кратчайшего двойственного сечения $J$ формула \eqref{eq:windows} и $\gamma_9=2$ дают $d(J)\le2D^{8/9}$. Точное сравнение
\[
2^9\cdot24450^8<15912^9
\]
показывает $d(J)<15912$. Поскольку в ранге 8 разрешены только остатки $0,1\pmod4$, получаем $d(J)\le15909$. Это ровно проверенное расширенное окно раздела 3; новая классификация за его пределами не предполагается.

Все прочие ветви имеют нижние границы выше 24450. Остаются $J_0,J_1,J_2$. При кратчайшем представителе клеящего класса последняя диагональ $n=s+\mu_J(a)$ является чётным числом не меньше 6. Для каждого из трёх родителей
\[
d(J)(8-R_J)>24450,
\]
так что $n=6$. Полное перечисление всех классов с
\[
s=6-\mu_J(a),\qquad D=d(J)s
\]
и проверкой \eqref{eq:extensionmin} в интервале $21870\le D\le24450$ даёт только $D=21870$: ровно прежние $1920+1+2$ класса. Все они изометричны $K_9$. Так как определитель ранга 9 чётен, вне этого класса следующий ещё не исключённый определитель равен 24452.
\end{proof}

Предложение не утверждает, что определитель 24452 реализуется. Оно устанавливает исключённый интервал, а не точное второе значение спектра определителей.

\begin{theorem}\label{thm:global10}
С учётом $\gamma_9=2$, исходного сертификата ранга 9 и новых результатов разделов 3--5 справедливо
\[
\boxed{31595\le e_{10}(6)\le54675}.
\]
\end{theorem}

\begin{proof}
Верхняя оценка следует из обеих матриц \eqref{eq:examples}. Если $D=d(L)\ge54675$, нижняя оценка очевидна. Иначе теорема \ref{thm:rank10e8} исключает наличие в $L$ любой копии $\sqrt3E_8$.

Кратчайшее двойственное сечение $J$ ранга 9 поэтому не изометрично $K_9$, поскольку $K_9$ содержит $\sqrt3E_8$. По предложению \ref{prop:gap} $d(J)\ge24452$. Теорема $\gamma_9=2$ и Морделл дают $\gamma_{10}\le2^{9/8}$. Следовательно,
\[
24452\le d(J)=D m(L^*)\le2^{9/8}D^{9/10},
\]
откуда
\[
D\ge24452^{10/9}/2^{5/4}=31594.9287\ldots.
\]
Округление проверено без вещественной арифметики: сравнение
\[
2^{45}D^{36}\ge24452^{40}
\]
впервые верно при целом $D=31595$. Это число также допустимо по сравнению $D\equiv0,3\pmod4$.
\end{proof}

С более слабыми внешними входами получаются более слабые, но полезные контрольные оценки. Одни классические константы до ранга 8 и исходное $e_9(6)$ дают $e_{10}(6)\ge24663$. Теорема Ноймайера $\gamma_{10}<16/7$ \cite{neumaier} даёт 26495. Использование $\gamma_9=2$ без нового разрыва даёт 27911. Число 31595 требует всей указанной цепочки новых результатов.

Следующий раздел усиливает эту предварительную оценку до точного равенства. Его новая полная цепочка не использует исходную нижнюю границу ранга 9 или прежний сертификат разрыва.

% Final insertion released after independent high- and low-range coverage audits.
% Final insertion released after independent high- and low-range coverage audits.
\section{Полный нижний предел в ранге 10}\label{sec:global10exact}

\begin{theorem}\label{thm:e10exact}
Всякая положительно определённая чётная целая решётка ранга 10 с квадратом минимума не меньше 6 имеет определитель не меньше 54675. Следовательно,
\[
\boxed{e_{10}(6)=54675}.
\]
Это компьютерно поддержанный результат с явно указанными ниже входами. Два известных примера $K_{10}$ и $\Kalt$ достигают границы; полная классификация \emph{глобальных} случаев равенства из данного строгого поиска не следует.
\end{theorem}

\subsection{Почему достаточно конечного расширенного окна}

Предположим, что существует контрпример $L_{10}$ с определителем $D<54675$. Чётность в ранге 10 даёт $D\equiv0,3\pmod4$, поэтому $D\le54672$. Теорема \ref{thm:rank10e8} исключает любую копию $\sqrt3E_8$ в $L_{10}$. Далее не используются ни прежняя нижняя граница для $e_9(6)$, ни предложение \ref{prop:gap}.

Публичная теорема $\gamma_9=2$ и неравенство Морделла дают $\gamma_{10}\le2^{9/8}$. Последовательное применение \eqref{eq:windows}, вместе с точными классическими константами до ранга 8 и сравнениями определителей, даёт:
\begin{center}
\begin{tabular}{crr}
\toprule
Ранг & Нижний предел & Верхний предел\\\midrule
1&6&6\\
2&27&44\\
3&108&210\\
4&324&789\\
5&972&2492\\
6&2187&6459\\
7&4374&13938\\
8&6564&24673\\
9&19684&40054\\
10&---&54672\\
\bottomrule
\end{tabular}
\end{center}
Все округления проверяются целыми степенями, без плавающей арифметики. Например, для самого верхнего перехода сравниваются
\[
d_9^{40}\le2^{45}D^{36},
\]
а для следующего $d_8^9\le2^9d_9^8$. Верхние вещественные целые пределы рангов 8 и 4 равны 24675 и 790; сравнение по модулю 4 уменьшает их до 24673 и 789.

Объясним нижние концы без обращения к старым вычислениям. В рангах 1--8 неравенство Эрмита даёт соответственно
\[
6,\ 27,\ 108,\ 324,\ 972,\ 2187,\ 4374,\ 6561.
\]
Равенство в ранге 8 возможно только для $\sqrt3E_8$, что уже исключено аналитически. Следующий разрешённый определитель равен 6564. В ранге 9 неравенство $\gamma_9\le2$ даёт $d_9\ge3^9=19683$; чётность определителя повышает предел до 19684. Для ранга 10 нижний конец не нужен. В частности, единственный начальный объект флага есть $[6]$; это вывод из окна, а не предположение о порождении минимальными векторами.

Полный интервал ранга 9 разбит на $[19684,24450]$ и $[24452,40054]$. Между ними пропущено только нечётное число 24451, невозможное в ранге 9. Граница разбиения не является входной теоремой о спектре определителей.

Любой предполагаемый контрпример имеет примитивный флаг кратчайших двойственных сечений целиком в этих окнах. При восстановлении флага снизу вверх достаточно перечислить все одношаговые расширения по \eqref{eq:extensiondet}--\eqref{eq:extensionmin}, требуя дополнительно, чтобы перпендикулярная двойственная нормаль была кратчайшей.

\subsection{Точный целочисленный тест двойственной нормали}

Опишем применённое ускорение, от которого зависит полнота поиска. Пусть родитель имеет матрицу $G$, определитель $d$, присоединённую матрицу $A=\operatorname{adj}G$ и новый вектор задаётся колонкой спариваний $p$ и чётной диагональю $n$. Определитель ребёнка равен
\[
D=nd-p^{\mathsf T}Ap,\qquad s=D/d.
\]
Норма двойственного вектора с координатами $(q,m)$ равна
\[
q^{\mathsf T}G^{-1}q+
\frac{(m-p^{\mathsf T}G^{-1}q)^2}{s}.
\]
При $q=0$ ненулевой вектор не короче перпендикулярной нормали нормы $1/s$. При $q\ne0$ минимизация по целому $m$ даёт расстояние до ближайшего целого. Поэтому кратчайшесть нормали равносильна точному целочисленному условию
\begin{equation}\label{eq:integerdualtest}
D(q^{\mathsf T}Aq)+
\operatorname{dist}(p^{\mathsf T}Aq,d\Z)^2\ge d^2
\end{equation}
для всех ненулевых $q$ с $D(q^{\mathsf T}Aq)<d^2$. Остальные $q$ заведомо безопасны.

Если $D\ge D_{\min}$, один полный список двойственных векторов родителя до квадрата нормы $d/D_{\min}$ содержит все необходимые $q$. Он строится прежним точным рациональным перечислителем; один представитель пары $\pm q$ достаточен, поскольку условие \eqref{eq:integerdualtest} не меняется при смене знака.

Класс с колонкой $p$ кодируется $Ap\bmod d$. Если $M_a$ --- целочисленный числитель его точного минимума, то условие слоёв записывается как
\begin{equation}\label{eq:integerlayertest}
m^2D+M_{ma}\ge6d
\qquad(m>0,\ m^2D<6d).
\end{equation}
Эта проверка включает ненулевые слои нулевого класса. Полные таблицы минимумов покрывают ровно $d$ классов; диагональ $n$ и сравнения для $D$ проверяются точно. Ускорение меняет форму вычисления, но не набор допустимых расширений.

Для полноты таблицы недостаточно увидеть по одному вектору каждого класса. Перечисляется весь двойственный шар до рациональной границы, при которой все $d$ классов уже покрыты; минимум в каждом классе тогда совпадает с глобальным минимумом этого класса. Полнота шара обеспечивается рекурсией по точному разложению $LDL^{\mathsf T}$: на каждом шаге все уже выделенные квадраты неотрицательны, а оставшийся бюджет даёт полный целочисленный интервал следующей координаты. Концы вычисляются рационально и целочисленным квадратным корнем, без округления вещественных чисел внутрь шара. Для каждого класса перебираются \emph{все} чётные $n$, попадающие в окно определителей; условие $n=6$ здесь не вводится.

У каждого используемого родителя отдельно проверяется равенство
\[
m(J^*)=
\frac{\det(\text{ведущего блока родителя})}{\det(\text{родителя})}.
\]
Таким образом, неравенство для следующего шага
\begin{equation}\label{eq:continuationfloor}
d_{r+1}\ge\frac{3d_r^2}{4d_{r-1}}
\end{equation}
применяется к подтверждённому кратчайшему двойственному флагу, а не к произвольно выбранному базисному сечению.

\subsection{Полное покрытие кандидатов и заключительное отсечение}

В расширенном окне получены следующие точные результаты. В столбце «осталось» указаны записи, не исключённые необходимой оценкой \eqref{eq:continuationfloor} для следующего ранга. Представители нужны для покрытия этих записей, а не для утверждения полной классификации всех решёток данного ранга.
\begin{center}
\begin{tabular}{crrrr}
\toprule
Ранг & Все записи & Отсечено для продолжения & Осталось & Представители\\\midrule
2&9&0&9&6\\
3&102&0&102&31\\
4&729&0&729&161\\
5&6934&0&6934&925\\
6&28284&0&28284&3752\\
7&56269&52208&4061&224\\
8&19481&18110&1371&47\\
9, низкий интервал&0&0&0&0\\
9, высокий интервал&3819&3819&0&0\\
\bottomrule
\end{tabular}
\end{center}

Для ранга 7 покрыты все 3752 родителя ранга 6. Из них только 500 допускают продолжение после предварительной необходимой оценки; остальные исключаются до дорогого перечисления. Для ранга 8 полностью обработаны 224 сохранённых представителя, а для каждого интервала ранга 9 --- все 47. Ни один родитель из требуемого входного списка не заменён выборочным пробным расчётом.

Каждая сохраняемая запись при сокращении до представителей снабжена явной унимодулярной изометрией; все такие матрицы проверены точным перемножением. Поэтому даже без предположения об идеальности алгоритма распознавания изометрий объединение классов не может удалить математически новый вариант без удостоверенного переноса на сохранённый представитель. Полнота исходного списка обеспечивается отдельно генератором расширений.

Изометрия не обязана сохранять прежний маркированный подфлаг: продолжение зависит от самой родительской решётки и всех её дискриминантных классов. Представитель выбран из фактически полученных матриц, поэтому он сам имеет корректный кратчайший двойственный флаг. Любое следующее расширение переносится изометрией на него. Оценка отсечения зависит от внутреннего инварианта $m(J^*)$, а не от произвольной маркировки; следовательно, другой выбор предыдущего флага не может восстановить отсечённый вариант.

Для высокого интервала ранга 9 у всех 3819 записей оценка \eqref{eq:continuationfloor} превосходит разрешённый верхний предел 54672. Минимальная из этих необходимых границ равна
\[
\frac{12966723}{212}=61163.7877\ldots>54672.
\]
Значит, ни одна из этих записей не может быть предпоследним членом флага предполагаемого контрпримера. Низкий интервал не дал ни одной записи; приведём для него дополнительно независимый короткий сертификат.

У 45 из 47 родителей ранга 8 явный двойственный вектор даёт по \eqref{eq:lift} необходимую нижнюю границу следующего определителя выше 24450. Для этой цели даже не требуется доказывать, что выбранный вектор кратчайший. У двух оставшихся родителей определители равны 14641 и 13689. Независимое точное перечисление двойственных векторов до квадрата нормы 4 покрывает соответственно все 14641 и 13689 дискриминантных классов. Поэтому у обоих $R_J\le4$. По \eqref{eq:extensionmin} слой $m=1$ требует $s\ge6-R_J\ge2$, и определитель продолжения не меньше соответственно 29282 или 27378. Оба числа больше 24450. Следовательно, весь низкий интервал исключён.

Таким образом, ни одна ветвь полного интервала ранга 9 не продолжается до контрпримера ранга 10. Непосредственное перечисление матриц ранга 10 уже не требуется: необходимое условие продолжения нарушается раньше.

Матрицы $K_{10}$ и $\Kalt$ имеют минимум 6 и определитель 54675, поэтому нижняя граница достигается. Это завершает доказательство теоремы \ref{thm:e10exact}.

Для одноточечной антиподной конструкции в размерности 38 при гарантированном радиусе $\sqrt6/2$ это даёт
\[
\delta_{38}\le\frac{(3/2)^{19}}{\sqrt{54675}}.
\]
Напечатанное вложение $K_{10}$ достигает этой величины. Следовательно, замена сечения ранга 10 с минимумом не меньше 6 не улучшает данную одноточечную конструкцию. Это не верхняя граница для всех упаковок размерности 38 и не решение общей задачи о многоточечных кластерах.

\subsection{Зависимости, воспроизводимость и граница классификации}

Полный вывод использует публичный результат $\gamma_9\le2$, классические константы Эрмита до ранга 8 и классификацию равенства в ранге 8, неравенство Морделла, дискриминантные минимумы $\sqrt3E_8$, теорему подкласса \ref{thm:rank10e8}, основанную на двумерной оценке Эрмита и новый полный флаговый перебор. \emph{Старый сертификат нижней границы ранга 9 и старое доказательство разрыва не являются зависимостями этого вывода.} Публичная классификация совершенных форм, лежащая за $\gamma_9=2$, не воспроизводится этим пакетом.

Оптимизированный генератор использует прежний рациональный перечислитель векторов. Независимая проверка полной цепочки охватывает точные окна, все родительские списки, отсечения и отображения на представителей; отдельным перечислителем перепроверены двойственные минимумы 8966 матриц. Низкий интервал дополнительно исключён программой, не импортирующей исходный генератор расширений или перечислитель векторов. Эти проверки являются внутренним аудитом, а не внешним рецензированием.

Кроме того, полный повторный запуск генерации заново воспроизвёл все 1897 требуемых списков родителей и 115627 исходных матриц; все результаты совпали с сохранёнными записями. Повторный запуск завершился статусом \file{PASS_FULL_EXACT_REPLAY}. Он использует проверенные генераторы и потому не заменяет независимую реализацию всей программы.

Основные файлы: \file{rank10/fast_extension_proof.txt},
\file{rank10/verify_global_rank10_certificate.py},
\file{independent_review/check_rank10_global_chain.py},
\file{independent_review/check_global_rank10_windows.py} и
\file{independent_review/check_rank9_low_by_radius.py}.

Установлено точное \emph{значение} глобального минимума. В теореме \ref{thm:rank10e8} ровно два случая равенства классифицированы только при наличии $\sqrt3E_8$. Поиск настоящего раздела использует строгое окно $D\le54672$; дополнительные решётки определителя 54675, не содержащие $\sqrt3E_8$, этим поиском не классифицированы. При равенстве верхние окна увеличиваются до $d_9\le40056$, $d_8\le24676$, $d_7\le13940$, $d_6\le6460$; эти дополнительные концы нельзя отбросить на основании большого запаса в последнем отсечении строгого поиска.

\subsection{Новое независимое следствие для ранга 9}

\begin{corollary}\label{cor:newrank9}
Публичная оценка $\gamma_9\le2$ вместе с новым флаговым перебором и таблицей классов $\sqrt3E_8$ даёт альтернативное доказательство:
\[
e_9(6)=21870,\qquad d(L)\le24450\Longrightarrow L\cong K_9.
\]
Оно не использует старый сертификат ранга 9.
\end{corollary}

\begin{proof}
Пусть $L$ имеет ранг 9, минимум не меньше 6 и $D\le24450$. Кратчайший двойственный родитель ранга 8 имеет определитель не больше 15909; весь его флаг лежит внутри нового, более широкого окна. Все применённые к этому префиксу необходимые фильтры остаются верными, поскольку последний определитель не больше 24450 и, следовательно, не больше 40054.

Если родитель не изометричен $\sqrt3E_8$, он попадает в новый список и исключается низким интервалом. Поэтому родитель есть $\sqrt3E_8$. Его класс имеет минимум $\mu\in\{0,2/3,4/3,2,8/3\}$. После выбора кратчайшего представителя присоединённый вектор имеет чётную норму $n\ge6$ и
\[
D=6561(n-\mu).
\]
При $n=6$ наименьшее значение равно 21870, следующее --- 26244. При $n\ge8$ имеем $D\ge34992$. Значит, при $D\le24450$ возможен только определитель 21870 и только один из 1920 глубоких классов с $\mu=8/3$. Их проверенная транзитивная орбита даёт одну изометрию $K_9$. Этот пример существует и имеет минимум 6, что доказывает оба утверждения.
\end{proof}

Таким образом, зависимости не образуют круга: новый перебор доказывает теорему \ref{thm:e10exact} непосредственно, а затем отдельно даёт следствие \ref{cor:newrank9}. Ранее изложенный путь к единственности сохраняет иной набор предпосылок и не требует публичной константы $\gamma_9=2$ в своём собственном случае равенства.

\subsection{Первые определители ранга 9 и первый случай без \texorpdfstring{$\sqrt3E_8$}{sqrt(3)E8}}

\begin{corollary}\label{cor:spectrum9}
При тех же предпосылках множество определителей чётных целых решёток ранга 9 с минимумом не меньше 6, строго меньших 32076, равно
\[
\{21870,\ 26244,\ 30618\}.
\]
При каждом из этих трёх определителей имеется ровно один класс изометрии. Таким образом, следующий определитель после $K_9$ равен точно 26244. Наименьший определитель решётки этого ранга и минимума, не содержащей $\sqrt3E_8$, равен точно 32076.
\end{corollary}

\begin{proof}
Наименьший определитель среди всех 3819 записей высокого интервала равен 32076. При выводе о ранге 9 не используем последнее отсечение, относящееся к рангу 10: нужен полный список \emph{до} него. Для любого $d_9<32076$ все необходимые верхние окна находятся внутри нового префикса, построенного для $d_9\le40054$. Поэтому кратчайший двойственный родитель обязан быть $\sqrt3E_8$.

Как в предыдущем доказательстве, $D=6561(n-\mu)$, $n\ge6$ чётно, а $\mu\in\{0,2/3,4/3,2,8/3\}$. При $n\ge8$ получаем $D\ge34992$. При $n=6$ ровно три указанных значения меньше 32076. Все достигаются: $s=6-\mu\ge10/3$ обеспечивает первый слой минимума 6, а слои $|m|\ge2$ имеют норму не меньше $4s>6$.

Точное вычисление отражений $E_8$ даёт транзитивность на каждом из трёх соответствующих множеств классов размеров 1920, 2240 и 2160. Эти целочисленные изометрии переносят одношаговые расширения и доказывают единственность при каждом определителе. Числа ориентированных минимальных векторов равны соответственно 258, 246 и 242.

В базисе $G_8$ раздела 4 примеры определителей 26244 и 30618 задаются матрицами $\left(\begin{smallmatrix}G_8&p\\p^{\mathsf T}&6\end{smallmatrix}\right)$ с колонками соответственно
\[
p=(1,-1,0,-1,-1,-3,0,0)^{\mathsf T},\qquad
p=(0,1,-1,1,0,1,2,2)^{\mathsf T}.
\]
Пример без $\sqrt3E_8$ имеет матрицу
\[
G_{32076}=
\begin{pmatrix}
6&-3&-3&-3&-3&-3&-3&-3&-3\\
-3&6&0&0&0&0&0&0&0\\
-3&0&6&0&0&3&0&0&1\\
-3&0&0&6&3&0&3&3&2\\
-3&0&0&3&6&0&3&3&2\\
-3&0&3&0&0&6&0&1&1\\
-3&0&0&3&3&0&6&3&3\\
-3&0&0&3&3&1&3&6&2\\
-3&0&1&2&2&1&3&2&6
\end{pmatrix}.
\]
Точная проверка даёт положительную определённость, определитель 32076, минимум 6 и 180 ориентированных минимальных векторов. Наличие $\sqrt3E_8$ означало бы её примитивность и представление
\[
D=3^7(3n-p^{\mathsf T}C^{-1}p),\qquad C=G_8/3.
\]
Поскольку $n$ чётно и $C^{-1}$ --- чётная целая матрица Грама, это число делится на $2\cdot3^7=4374$. Но $32076\equiv1458\pmod{4374}$, так что данный пример не содержит $\sqrt3E_8$. Нижняя граница уже доказана полнотой префикса.
\end{proof}

Точные примеры и проверки сохранены в \file{rank10/rank9_low_spectrum.json} и \file{rank10/verify_rank9_low_spectrum.py}. Вывод об отсутствии меньших определителей использует полный сертификат префикса; проверка одной матрицы его не заменяет. Единственность при определителе 32076 не утверждается.


\section{Кластерные расширения: общая геометрия}

\subsection{Как сравнивается плотность}

Пусть $K$ --- примитивное сечение ранга $k$ чётной унимодулярной решётки ранга 48 с минимумом 6. Для конечного кластера $S\subset K^*$ положим
\[
s=|S|,\qquad\beta=\max_{u,v\in S}\|u-v\|^2<6.
\]
Антиподная конструкция гарантирует радиус шаров $\sqrt{6-\beta}/2$ в размерности $n=48-k$ и центральную плотность
\begin{equation}\label{eq:density}
\delta(S,K)=\frac{s(6-\beta)^{n/2}}{2^n\sqrt{d(K)}}.
\end{equation}
Если конкретное окружающее вложение даёт больший фактический минимум, его улучшенный радиус требует отдельного анализа. Все утверждения этого отчёта относятся к радиусу, гарантированному именно $6-\beta$.

Относительно одноточечного кластера над тем же $K$ коэффициент равен
\begin{equation}\label{eq:ratio}
R=s(1-\beta/6)^{n/2}.
\end{equation}
Для $\beta<6$ разные точки $S$ принадлежат разным классам $K^*/K$, иначе их разность --- ненулевой вектор $K$ нормы меньше 6. Поэтому $s\le d(K)$.

\subsection{Проекция на восемь координат и её слои}

Пусть $J\cong\sqrt3E_8\subset K$. Ортогональная проекция $K^*$ на $\spanR J$ лежит в $J^*$: это следует из целочисленных спариваний с $J$. Поэтому все ненулевые квадраты расстояний между проекциями принадлежат $(2/3)\Z$.

Каждый непустой слой проекции является сдвигом решётки
\[
T=K^*\cap(\spanR J)^\perp.
\]
Если матрица $K$ записана блоками $\begin{psmallmatrix}A&B\\B^{\mathsf T}&D\end{psmallmatrix}$, то матрица $T$ равна обратной матрице дополнения Шура $D-B^{\mathsf T}A^{-1}B$. Для наших примеров:
\begin{equation}\label{eq:fibers}
K_9:\ T=[3/10],\qquad
K_{10},\Kalt:\ T=\frac15\begin{pmatrix}2&-1\\-1&2\end{pmatrix}\cong A_2/\sqrt5.
\end{equation}
Ни сюръективность проекции, ни её конкретный клеящий характер в последующих оценках не требуются.

\subsection{Число точек с малым числом расстояний}

Эквидистантное множество в $\R^8$ имеет не больше 9 точек. Общая удобная оценка для множества $P\subset\R^8$ с $r$ различными положительными квадратами расстояний такова:
\begin{equation}\label{eq:fewdist}
|P|\le\binom{9+r}{r}.
\end{equation}
Для каждого $x\in P$ рассмотрим
\[
f_x(y)=\prod_{i=1}^r(\|x-y\|^2-d_i).
\]
Матрица значений на $P$ диагональна с ненулевой диагональю, поэтому $f_x$ линейно независимы. Разложение $\|x-y\|^2=\|y\|^2-2(x,y)+\|x\|^2$ помещает их в сумму пространств
\[
\sum_{j=0}^r\|y\|^{2j}\mathcal P_{\le r-j}(\R^8).
\]
Сумма размерностей равна $\sum_{k=0}^r\binom{8+k}{k}=\binom{9+r}{r}$, что и доказывает \eqref{eq:fewdist}. Оценка не оптимальна, но достаточна.

Для проекции кластера отсюда следуют пределы 9 при $\beta<4/3$, 55 при $\beta<2$, 220 при $\beta<8/3$, 715 при $\beta<3$ и 24310 при $\beta<6$. Первый предел использует более сильную эквидистантную оценку.

\section{Отсутствие улучшения во всём семействе ранга 9}

\begin{theorem}\label{thm:cluster9}
Пусть $K$ --- любая чётная целая решётка ранга 9 с $m(K)\ge6$, содержащая $\sqrt3E_8$, и $S\subset K^*$ --- конечный кластер с $\beta<6$. Тогда
\[
\delta(S,K)\le\frac{6^{39/2}}{2^{39}\sqrt{21870}}.
\]
Равенство возможно только для $s=1$ и $d(K)=21870$. Для всякого нетривиального кластера отношение к правой части не превосходит
\[
9(8/9)^{39/2}=0.9052495585\ldots<1.
\]
Оценка 0.90525 не утверждается точной.
\end{theorem}

\begin{proof}
Примитивность вложенного $J=\sqrt3E_8$ автоматическая, как в разделе 4. Из \eqref{eq:gluebound} следует $D=d(K)\ge21870$. Одномерный слой $T$ имеет шаг
\[
\tau=\sqrt{6561/D}\le\tau_0=\sqrt{3/10}.
\]
Введём $f(\beta)=(1-\beta/6)^{39/2}$. Отношение к одноточечной упаковке над $K_9$ равно
\begin{equation}\label{eq:familyratio}
R=\sqrt{10/3}\,\tau s f(\beta).
\end{equation}
Пусть $M$ --- число непустых проекционных слоёв, а $s_i$ --- число точек в $i$-м слое. Их вертикальный размах не меньше
\[
\ell_i=(s_i-1)\tau,\qquad \ell_i\le\sqrt\beta.
\]
Максимальное вертикальное расстояние между двумя слоями не меньше $(\ell_i+\ell_j)/2$. Поэтому для разных проекций $p_i,p_j$
\begin{equation}\label{eq:crossspans}
\beta\ge\|p_i-p_j\|^2+(\ell_i+\ell_j)^2/4.
\end{equation}

\textbf{Одна проекция.} При $M=1$ и $s\ge2$ имеем $\beta\ge(s-1)^2\tau^2$, откуда
\[
R\le\frac{s}{s-1}\sqrt{10\beta/3}\,f(\beta).
\]
Максимум $\sqrt\beta f(\beta)$ достигается при $\beta=3/20$. Следовательно,
\[
R\le\sqrt2(39/40)^{39/2}<1.
\]
Для $s=1$ имеем $\beta=0$ и $R=\sqrt{21870/D}\le1$, с указанным случаем равенства.

\textbf{Несколько проекций, $\beta<4/3$.} Необходимо $\beta\ge2/3$, а $M\le9$. Если все слои одноточечные, то
\[
R\le9f(2/3)=9(8/9)^{39/2}<1.
\]
Иначе обозначим через $t\ge u$ два наибольших числа точек в слоях. Тогда $t\ge2$, $u\ge1$, $s\le t+(M-1)u$ и по \eqref{eq:crossspans}
\[
\tau(t+u-2)\le2\sqrt{\beta-2/3}.
\]
Поскольку
\[
\frac{t+(M-1)u}{t+u-2}\le M+1
\]
(разность после умножения равна $M(t-2)+2(u-1)\ge0$), получаем
\[
R\le20\sqrt{(10/3)(\beta-2/3)}\,f(\beta).
\]
Максимум правой части достигается при $\beta=4/5$ и равен
\[
\frac{40}{3}(13/15)^{39/2}<1.
\]
Оба вычисления критических точек следуют из
\[
\frac{d}{d\beta}\log\big((\beta-c)^{1/2}f(\beta)\big)=0
\Longleftrightarrow 6-\beta=39(\beta-c).
\]

\textbf{Промежуток $4/3\le\beta<2$.} Назовём слой тяжёлым, если $\ell_i>\sqrt{\beta-4/3}$. Проекции двух тяжёлых слоёв не могут находиться на квадрате расстояния $4/3$, иначе \eqref{eq:crossspans} нарушается. Их расстояние обязательно $2/3$, поэтому тяжёлых слоёв не больше 9. Всего слоёв не больше 55, откуда
\[
\tau s=M\tau+\sum_i\ell_i
\le55\tau_0+46\sqrt{\beta-4/3}+9\sqrt\beta.
\]
При $4/3\le\beta<3/2$, используя $\sqrt{5/9}<3/4$ и $\sqrt5<9/4$, выводим
\[
R<\frac{439}{4}f(4/3)<1.
\]
При $3/2\le\beta<2$ оценки $\sqrt{20/9}<3/2$, $\sqrt{20/3}<8/3$ дают
\[
R<148f(3/2)<1.
\]

\textbf{Оставшиеся диаметры.} Всегда $s_i\le1+\sqrt\beta/\tau$, поэтому
\[
R\le M(1+\sqrt{10\beta/3})f(\beta).
\]
С использованием \eqref{eq:fewdist} и рациональных верхних оценок квадратных корней получаем таблицу:
\begin{center}
\begin{tabular}{crr}
\toprule
Интервал $\beta$ & Предел для $M$ & Верхняя оценка $R$\\\midrule
$[2,8/3)$ &220&$880f(2)$\\
$[8/3,3)$ &715&$(12155/4)f(8/3)$\\
$[3,6)$ &24310&$133705f(3)$\\
\bottomrule
\end{tabular}
\end{center}
Каждое число строго меньше 1. Все сравнения выполнены точно после возведения в квадрат; например,
\[
2(39/40)^{39}<1,\qquad (1600/9)(13/15)^{39}<1.
\]
Так же проверено, что все полученные верхние границы не больше $9(8/9)^{39/2}$. Случаи покрывают все $\beta<6$, что завершает доказательство.
\end{proof}

Это утверждение охватывает бесконечное семейство больших определителей. Оно не основано на ограниченном переборе решёток или кластеров. Однако отсутствие $\sqrt3E_8$ в произвольной решётке ранга 9 выводит её за область доказательства.

\section{Все диаметры для обоих минимальных типов ранга 10}

\begin{theorem}\label{thm:cluster10}
Для $K=K_{10}$ или $K=\Kalt$, всякого конечного $S\subset K^*$ с $s>1$ и $\beta<6$ имеем
\[
s(1-\beta/6)^{19}<1.
\]
Значит, одноточечный кластер строго оптимален при гарантированном радиусе антиподной конструкции. В силу теоремы \ref{thm:rank10e8} это охватывает все минимизаторы подкласса ранга 10, содержащего $\sqrt3E_8$.
\end{theorem}

\subsection{Конечные задачи в треугольном слое}

Слой \eqref{eq:fibers} имеет квадратичную форму
\[
\| (u,v)\|_T^2=\frac25(u^2-uv+v^2).
\]
Положительные нормы ниже $14/5$ равны $2/5,6/5,8/5$. Максимальные размеры кластеров в слое таковы:
\begin{center}
\begin{tabular}{crrr}
\toprule
Интервал квадрата диаметра &$[2/5,6/5)$&$[6/5,8/5)$&$[8/5,14/5)$\\\midrule
Максимальный размер &3&4&7\\
\bottomrule
\end{tabular}
\end{center}
Здесь использованы небольшие полные конечные сертификаты. После переноса одной точки в нуль соответствующие шары содержат 7, 13 и 19 точек. Для опровержения границ требуется соответственно выбрать 3, 4 и 7 других точек. Проверены все
\[
\binom63=20,\qquad \binom{12}4=495,\qquad \binom{18}7=31824
\]
подмножества. Допустимых множеств запрещённого размера нет; примеры размеров 3, 4, 7 сохранены. Полнота шаров следует из
\[
2(u^2-uv+v^2)\ge\frac32u^2,\qquad
2(u^2-uv+v^2)\ge\frac32v^2:
\]
при норме меньше $14/5$ достаточно проверить все 49 пар $|u|,|v|\le3$. Все расстояния сравниваются как целые значения $2(u^2-uv+v^2)$.

\subsection{Несколько многоточечных слоёв}

Две точки в одном слое и точка с другой проекцией вынуждают
\[
\beta\ge2/3+(2/5)/4=23/30.
\]
Это обычная оценка через середину отрезка: хотя бы один конец отрезка длины не меньше $\sqrt{2/5}$ удалён от любой точки вертикально не меньше чем на половину длины.

Для двух многоточечных слоёв возьмём ненулевые внутренние разности $a,b\in T$. Четыре поперечных разности имеют вид $z+(\pm a\pm b)/2$. Сравнение противоположных пар показывает, что максимальная вертикальная норма не меньше
\[
\frac14\max(\|a+b\|^2,\|a-b\|^2)\ge3/10.
\]
Если один из векторов не минимален, то $(\|a\|^2+\|b\|^2)/4\ge2/5$. Если оба минимальны, они принадлежат шести корням треугольной решётки и $|(a,b)|\ge1/5$, что даёт $3/10$. Все 36 пар корней дополнительно проверены.

Таким образом, при проекциях $p,q$ двух многоточечных слоёв
\begin{equation}\label{eq:twoheavy10}
\beta\ge\|p-q\|^2+3/10.
\end{equation}
При $\beta<29/30$ многоточечный слой не более одного. При $\beta<49/30$ проекции всех многоточечных слоёв эквидистантны с квадратом расстояния $2/3$, так что их не больше 9, даже если все проекции допускают два расстояния.

\subsection{Таблица, покрывающая все диаметры}

Объединяя пределы проекций и слоёв, получаем следующие верхние границы $B$ для $s$. Все интервалы замкнуты слева и открыты справа.
\begin{longtable}{P{0.22\linewidth}rP{0.59\linewidth}}
\toprule
Интервал $\beta$ &$B$&Обоснование\\\midrule
\endhead
$[2/5,2/3)$&3&Одна проекция.\\
$[2/3,23/30)$&9&Либо все слои одноточечные, либо проекция одна.\\
$[23/30,29/30)$&11&Не больше 9 проекций; один слой может иметь 3 точки.\\
$[29/30,6/5)$&27&Не больше 9 слоёв по 3 точки.\\
$[6/5,4/3)$&36&Не больше 9 слоёв по 4 точки.\\
$[4/3,8/5)$&82&Не больше 55 слоёв, из них не больше 9 многоточечных по 4 точки: $55+9\cdot3$.\\
$[8/5,49/30)$&109&Не больше 55 слоёв, из них не больше 9 многоточечных по 7 точек: $55+9\cdot6$.\\
$[49/30,2)$&385&Не больше 55 слоёв по 7 точек.\\
$[2,8/3)$&1540&Не больше 220 слоёв по 7 точек.\\
$[8/3,6)$&54675&Инъекция кластера в $K^*/K$.\\
\bottomrule
\end{longtable}
Нетривиальный кластер с $\beta<2/5$ невозможен. Для каждой строки с левым концом $b$ точная арифметика даёт
\[
B^2(1-b/6)^{38}<1.
\]
Функция $(1-\beta/6)^{19}$ убывает, поэтому нужное строгое неравенство справедливо на всём интервале. Это доказывает теорему \ref{thm:cluster10} без ограничения числа точек, диаметра 1 или размера графа поиска.

Для $K_9$ также выполнен отдельный более короткий табличный сертификат с одномерным слоем шага $\sqrt{3/10}$. Он теперь следует и из более сильной теоремы \ref{thm:cluster9}; его независимое наличие служит дополнительной проверкой фиксированного случая.

\section{Что проверено для окружающих решёток ранга 48}

\subsection{Матрицы исходных вложений}

Из исходного PDF извлечены матрицы окружающих решёток $P_{48p}$ и $P_{48m}$ и матрицы их выбранных сечений. При соглашении, что строки $T$ задают базис сечения в окружающих координатах, проверены
\[
G_K=T G_{48}T^{\mathsf T},\qquad
T G_{48}B^{\mathsf T}=0,
\]
где строки $B$ задают полный целый базис ортогонального дополнения. Проверены симметрия, чётность, положительная определённость и определитель 1 обеих окружающих матриц; примитивность $T$ удостоверена нормальной формой Смита с ненулевыми инвариантами 1.

\begin{center}
\begin{tabular}{lrrl}
\toprule
Сечение &Ранг дополнения&Определитель&Инварианты $A_K$\\\midrule
$K_8\subset P_{48p}$&40&6561&$3,3,3,3,3,3,3,3$\\
$K_9\subset P_{48p}$&39&21870&$3,3,3,3,3,3,30$\\
$K_{10}\subset P_{48p}$&38&54675&$3,3,3,3,15,45$\\
$K_5\subset P_{48m}$&43&3072&$4,16,48$\\
$K_3\subset P_{48m}$&45&128&$2,8,8$\\
\bottomrule
\end{tabular}
\end{center}

Для каждого дополнения проверен тот же определитель и сохранён явный вектор нормы 6. Утверждение о минимуме 6 \emph{всей} окружающей решётки остаётся входом её известной экстремальности; новая проверка положительности и унимодулярности сама по себе не исключает нормы 2 и 4 в размерности 48.

При примитивном $K\subset\Lambda$ с унимодулярной $\Lambda$ стандартное точное отображение $\Lambda\to K^*$ сюръективно, а ядро равно $K^\perp\cap\Lambda$. Отсюда
\[
\Lambda/(K\oplus K^\perp)\cong K^*/K,\qquad
d(K^\perp)=d(K).
\]
В чётном случае индуцированное соответствие дискриминантных групп меняет знак их квадратичных форм. Такие соответствия для напечатанных матриц также проверены.

\subsection{Невложимость альтернативного типа над фиксированным сечением}

\begin{theorem}\label{thm:nonembedding}
Никакая изометричная копия $\Kalt$ в напечатанной решётке $P_{48p}$ не содержит конкретное напечатанное сечение $K_9\subset P_{48p}$.
\end{theorem}

Зафиксируем \emph{конкретный маркированный} базис $t_1,\ldots,t_9$ напечатанного $K_9\subset P_{48p}$. Исследована система
\begin{equation}\label{eq:marked}
\|w\|^2=6,\qquad ((t_i,w))_{i=1}^9=(-p,-1).
\end{equation}
Если бы она имела решение $w\in P_{48p}$, добавление $w$ давало бы именно блочную матрицу $\Kalt$ из \eqref{eq:examples} с этим маркированным $K_9$.

Полный точный поиск решения \eqref{eq:marked} свёл его к задаче ближайшего вектора в аффинном классе ортогонального дополнения ранга 39. Независимая реализация Python с неограниченными целыми и рациональными числами прошла 611239 узлов и не нашла ни одного допустимого листа. Результат совпал с реализацией C++. Никакого вещественного отсечения при итоговой сертификации не применялось.

В качестве положительного контроля тот же алгоритм выполнен для колонки скалярных произведений исходного $K_{10}$. Он прошёл 610015 узлов и нашёл ровно один вектор, совпавший с последней строкой напечатанной матрицы вложения $K_{10}$.

\begin{proof}[Почему пустой маркированный слой исключает любую копию над данным $K_9$]
Предположим, что некоторая копия $\Kalt$ содержит фиксированное $K_9$, а значит, и его фиксированное $J=\sqrt3E_8$. Проекция этой копии на двумерное ортогональное дополнение к $J$ имеет минимум не меньше $10/3$ и определитель $25/3$. Как в теореме \ref{thm:rank10e8}, она обязана быть шестиугольной. Проекция $z_1$ девятого фиксированного вектора имеет норму $10/3$, так что её можно дополнить остроугольным базисом $z_1,z_2$ с матрицей \eqref{eq:binaryprojection}. Выберем кратчайший подъём $w$ второго вектора; он имеет полную норму 6, а его $J$-компонента --- норму $8/3$.

Пусть $q'$ --- колонка скалярных произведений $w$ с фиксированным базисом $J$. Если $q'$ и $p$ независимы по модулю 3, то матрица
\[
\begin{pmatrix}0&p&q'\\p^{\mathsf T}&0&c\\(q')^{\mathsf T}&c&0\end{pmatrix}
\]
имеет ранг 4 над $\F_3$. Это противоречит инварианту $\rank_{\F_3}\Kalt=3$. Поэтому $q'\equiv0,p$ или $-p\pmod3$. Нулевой класс невозможен: его минимум 0, а подъём при минимальной проекции требует $8/3$. Класс $p$ невозможен из целочисленности полного спаривания: $J$-часть даёт $8/3\pmod\Z$, перпендикулярная часть $5/3$, и сумма $13/3$ не целая.

Значит, $q'\equiv-p\pmod3$. Так как $C$ унимодулярна, это равносильно совпадению дискриминантного класса с отрицанием класса девятого вектора. Вычитанием вектора $J$ можно сделать $J$-проекцию $w$ точно равной отрицательной $J$-проекции $t_9$. Новая проекция имеет ту же минимальную норму $8/3$, поэтому полная норма остаётся 6, а спаривание с $t_9$ равно
\[
-8/3+5/3=-1.
\]
Получен вектор системы \eqref{eq:marked}, противоречащий её точной пустоте. Это доказывает теорему, не требуя классификации базисов или копий $E_8$ внутри $K_9$.
\end{proof}

\subsection{Почему целочисленный поиск действительно полон}

Положим $A=T_9G_{48}$, $b=(-p,-1)$. Целый правый обратный оператор $Q$ даёт частное решение $w_0=Qb$, а проверенный целый базис ядра $B$ даёт все решения линейных условий:
\[
w=w_0+B^{\mathsf T}x,\qquad x\in\Z^{39}.
\]
Обозначим $C_0=BG_{48}B^{\mathsf T}$, $g=BG_{48}w_0$, $c=-C_0^{-1}g$. Точно вычислено
\[
w_0^{\mathsf T}G_{48}w_0-g^{\mathsf T}C_0^{-1}g=7/2.
\]
Поэтому условие нормы не больше 6 равносильно
\[
(x-c)^{\mathsf T}C_0(x-c)\le5/2.
\]
Точное разложение $C_0=L\diag(D_i)L^{\mathsf T}$ превращает левую часть в
\[
\sum_iD_i\left(x_i+\sum_{j>i}L_{ji}x_j-y_i\right)^2,\qquad y=L^{\mathsf T}c.
\]
Из неравенства Коши получены доказанные целые границы $|x_i|\le B_i$, удовлетворяющие
\[
B_i\ge|c_i|,\qquad (B_i-|c_i|)^2\ge(5/2)(C_0^{-1})_{ii}.
\]
Максимальная граница $B_i$ равна 28.

Для ускорения используется \emph{внешняя целочисленная релаксация}, а не доверие к вещественному округлению. При $S=2^{24}$ округляем $SL_{ji}$ и $Sy_i$ до целых с ошибкой не больше $1/2$, кладём $d_i=\lfloor SD_i\rfloor>0$ и для зафиксированного хвоста определяем
\[
C_i=\operatorname{round}(Sy_i)-\sum_{j>i}x_j\operatorname{round}(SL_{ji}),\qquad
E_i=\left\lceil\frac{1+\sum_{j>i}B_j}{2}\right\rceil.
\]
Точный центр отличается от $C_i/S$ не больше чем на $E_i/S$. Поэтому каждый точный вклад не меньше
\[
\frac{d_i}{S^3}\max(0,|Sx_i-C_i|-E_i)^2.
\]
Всякое подлинное решение обязано лежать в релаксации
\[
\sum_i d_i\max(0,|Sx_i-C_i|-E_i)^2\le5S^3/2.
\]
При оставшемся целом бюджете $R$ точный охватывающий интервал следующей координаты задаётся
\[
|Sx_i-C_i|\le E_i+\operatorname{isqrt}(\lfloor R/d_i\rfloor).
\]
Перебор всех этих интервалов с доказанными $B_i$ не даёт ни одного листа даже в расширенном множестве, а значит, исходная система пуста. Независимая Python-проверка заново удостоверяет ядро, разложение, все координатные и округлительные оценки; подготовленные коэффициенты не принимаются на доверии.

Теорема \ref{thm:nonembedding} \textbf{не} исключает все вложения $\Kalt\hookrightarrow P_{48p}$. Копия, содержащая иное сечение $K_9$, может относиться к другой орбите окружающей группы автоморфизмов. Эти орбиты полностью не перечислены. Проверка исходных матриц и новый отрицательный результат прошли дополнительный внутренний аудит, но не заменяют полной классификации.

\subsection{Все 19 минимальных надсечений фиксированного \texorpdfstring{$K_9$}{K9}}

\begin{theorem}\label{thm:nineteen}
В напечатанной решётке $\Lambda=P_{48p}$ имеется ровно 19 различных примитивных подрешёток $S$ ранга 10, содержащих фиксированное напечатанное $K_9$, с $d(S)=54675$ и $m(S)\ge6$. Все 19 изометричны исходному $K_{10}$. Это минимальный возможный определитель для надсечений данного $K_9$.
\end{theorem}

Число 19 считает \emph{конкретные подрешётки в окружающих координатах}, а не орбиты группы автоморфизмов. Оно относится к одному фиксированному сечению $K_9$, а не ко всем вложениям абстрактной решётки.

\begin{proof}
Каждое такое $S$ содержит фиксированное $\sqrt3E_8$, поэтому теорема \ref{thm:rank10e8} даёт нижнюю границу 54675 и шестиугольную проекцию при равенстве. Пусть $a=C^{-1}p\pmod3$ --- класс девятого фиксированного вектора. Три класса $a,b,a-b$ остроугольного проекционного базиса глубокие. Ровно 561 класс $b$ удовлетворяет этому условию. Инволюция
\[
b\longmapsto a-b
\]
имеет единственную неподвижную точку $b=-a$ и 280 двухэлементных орбит. Значит, достаточно 281 неориентированного целевого класса.

Выберем представителя $v$ нормы 8 в каждом классе $b$ и положим
\[
q=Cv,\qquad t=(p^{\mathsf T}v+5)/3,\qquad r_b=(q,t)^{\mathsf T}.
\]
Целочисленность $t$ следует из глубины $a-b$. Добавлением вектора $J$ к кратчайшему подъёму можно сделать его $J$-компоненту точно равной $v/3$. Поэтому всякое искомое надсечение появляется в одной из систем
\[
T_9G_{48}u=r_b,\qquad u^{\mathsf T}G_{48}u=6,\qquad u\in\Z^{48}.
\]
Прямая проверка даёт $r_b^{\mathsf T}K_9^{-1}r_b=7/2$ для всех целей. Следовательно, каждый поиск имеет остаточную норму $5/2$ в том же полном целочисленном ядре ранга 39 и подчиняется доказанной выше внешней целочисленной релаксации.

Полное перечисление всех 281 целей даёт:
\begin{center}
\begin{tabular}{lr}
\toprule
Точные поисковые узлы &171472428\\
Пустые целевые слои &262\\
Слои с единственным вектором нормы 6 &19\\
Всего листьев и решений &19\\
\bottomrule
\end{tabular}
\end{center}
Одноточечная орбита $b=-a$ пуста; все решения относятся к исходному типу $K_{10}$. Полная Python-проверка с неограниченными целыми повторяет рациональные предпосылки и каждый целевой поиск, сравнивая число узлов и список векторов с реализацией C++. Отдельный построитель из задачи ранга 10 даёт тот же набор из 561 дискриминантных классов; равенство наборов проверено по точным отпечаткам через присоединённую матрицу.

Остаётся удостоверить, что считаются подрешётки, а не несколько генераторов одной. Выбор представителя пары $b,a-b$ устраняет две остроугольные ориентации проекционного базиса. Для остальных 560 ориентированных классов порядок в $K_9^*/K_9$ равен 6; противоположные генераторы не могут иметь одну и ту же $K_9$-проекцию, иначе удвоенная проекция принадлежала бы $K_9$. Генераторы одинакового знака с той же проекцией также совпадают.

Дополнительно выполнена прямая проверка в окружающих координатах. Для каждого найденного $u$ строки $T_{10}=[T_9;u]$ имеют форму Смита $[I_{10}\mid0]$. Поэтому их решётка насыщена и является полным сечением своего вещественного пространства. При $d=21870$ введём целый вектор
\[
N(u)=d\,u-T_9^{\mathsf T}\operatorname{adj}(K_9)r_b.
\]
$N(u)/d$ есть перпендикулярная к фиксированному $K_9$ компонента и имеет норму $5/2$. Два пространства надсечений совпадают тогда и только тогда, когда их $N$ отличаются знаком: все ненулевые нормали имеют одинаковую длину. Все 19 векторов $N$ попарно различны с точностью до знака. Следовательно, получены ровно 19 разных насыщенных подрешёток. Полные координаты векторов, матрицы Грама, формы Смита и нормали сохранены в каталоге.
\end{proof}

Эта классификация завершает задачу оптимальных надсечений \emph{данного} вложения $K_9$. Она не определяет, сколько из 19 сечений эквивалентны под стабилизатором $K_9$ в окружающей группе автоморфизмов, и не классифицирует остальные орбиты вложений $K_9$.

Теоремы \ref{thm:nonembedding} и \ref{thm:nineteen} не используют глобальную нижнюю оценку $e_9(6)$ или её исходный большой перебор. Их входы --- точные напечатанные матрицы и примитивность, конечная таблица классов $E_8/3E_8$, элементарная двумерная оценка и полный поиск в указанных аффинных слоях. При явном условии $m(S)\ge6$ также не требуется заново доказывать минимум 6 всей окружающей решётки; допустимость найденных внутренних матриц проверена независимо.

\section{Литература, проверки и оставшиеся задачи}

\subsection{Что выяснено о научной новизне}

В работе Dutour Sikirić--van Woerden \cite{dw}, версия 2 от 7 февраля 2026 года, приведена классификация 2237251040 совершенных классов подобия ранга 9 и вывод $\gamma_9=2$. В ней имеется известный совершенный класс $Q_{129}$ с 258 минимальными векторами, связанный с более ранними работами Watson и описанный также в \cite{vwthesis}.

Для нашего $K_9$ теперь имеется точный сертификат совершенства: 258 минимальных векторов дают 129 противоположных пар, а их симметрические внешние произведения порождают пространство размерности
\[
\dim\operatorname{Sym}_9=45.
\]
Указан конкретный минор размера $45\times45$ с определителем 3. Поэтому $K_9$ является совершенной формой и относится к уже известному классу подобия $Q_{129}$ с этим числом минимальных векторов.

Следовательно, нельзя представлять саму решётку $K_9$ как впервые открытую в этой ветке. Предмет возможной новой статьи --- глобальная \emph{чётная целочисленная} задача минимума 6, полнота её случая равенства и дополнительные результаты данного отчёта. Совершенство одного примера не доказывает глобальную дискретную оптимальность: целочисленный минимизатор не обязан автоматически быть совершенной вещественной формой.

Проверены первичные источники по интегральным ламинированным решёткам \cite{cs,pp}, таблицы Martinet--Batut \cite{martinet} и классификация сильно совершенных решёток ранга 10 \cite{nv}. Их предпосылки отличаются: рекурсивная ламинация, примеры из каталога или сильное совершенство не эквивалентны произвольной чётной целой решётке ранга 10 с минимумом 6. В просмотренных утверждениях полного ответа на нашу глобальную задачу не найдено. Само отсутствие найденной ссылки не является доказательством новизны.

Неопубликованная рукопись, её название и новые целевые определители не использовались как поисковые запросы внешним сервисам. Коллегам и рецензентам сообщения в рамках текущего запуска не отправлялись.

\Needspace{12cm}
\subsection{Карта воспроизводимости}

Все пути ниже относительны к каталогу \file{research_20261007}. Использовать обычный Python 3 с включёнными \texttt{assert}; режим \texttt{-O} недопустим. Исходный точный пакет находится в подкаталоге \file{recovered/}.

\begin{longtable}{P{0.34\linewidth}P{0.59\linewidth}}
\toprule
Утверждение & Главный проверяющий файл\\\midrule
\endhead
Новые окна и уникальность $K_9$ &\file{equality/verify_equality.py}\\[3pt]
Разрыв до определителя 24450 &\file{equality/rank9_gap.py}\\[3pt]
Независимые граничные сертификаты и изометрии &\file{independent_review/check_equality_independent.py}\\[3pt]
Совершенство $K_9$ &\file{independent_review/check_k9_perfect.py}\\[3pt]
Подкласс ранга 10 и два типа &\file{rank10/verify_rank10_e8_classification.py}\\[3pt]
Предварительная нижняя граница 31595 &\file{rank10_theory/verify_rank10_bounds.py}\\[3pt]
Точное глобальное $e_{10}(6)$ &\file{rank10/verify_global_rank10_certificate.py}\\[3pt]
Независимая полная цепочка &\file{independent_review/check_rank10_global_chain.py}\\[3pt]
Независимое низкое окно ранга 9 &\file{independent_review/check_rank9_low_by_radius.py}\\[3pt]
Первые три определителя и случай без $\sqrt3E_8$ &\file{rank10/verify_rank9_low_spectrum.py}\\[3pt]
Равномерные частичные кластеры ранга 10 &\file{clusters/verify_rank10_single_projection.py}\\[3pt]
Фиксированные кластеры всех диаметров &\file{clusters/verify_fixed_section_cluster_optimality.py}\\[3pt]
Все сечения ранга 9 с $\sqrt3E_8$ &\file{clusters/verify_all_E8_rank9_cluster_optimality.py}\\[3pt]
Напечатанные вложения и дополнения &\file{embeddings/verify_ambient_embeddings.py}\\[3pt]
Маркированное продолжение и контроль &\file{embeddings/verify_A10_extension_exact.py}\\
Все 19 надсечений фиксированного $K_9$ &\file{embeddings/verify_all_fixed_K9_extensions.py}\\
\bottomrule
\end{longtable}

Таблицы JSON и журналы расположены рядом с программами. Контрольные суммы удостоверяют сохранность файлов; полноту доказывают точные окна, порядки групп, полные шары с рациональными границами, орбитные деревья и унимодулярные изометрии.

Исходный перебор и исправления разобраны в \cite{audit}. Новые окна и орбиты вычислены заново. Независимая проверка 262 границ использует иной метод; изометрии проверены повторным точным перемножением. В орбитном перечислении ранга 10 использован прежний арифметический модуль. Полной независимой второй реализации первоначального доказательства нет.

Кластерные программы проверяют малые конечные слои и рациональные неравенства; геометрические аргументы изложены выше. Внутренние рецензирующие заметки по новым результатам включены в пакет.

\subsection{Точный оставшийся объём работы}

\begin{enumerate}
\item \textbf{Глобальное равенство в ранге 10.} Значение $e_{10}(6)=54675$ установлено; требуется классифицировать все решётки этого определителя без предположения о наличии $\sqrt3E_8$. Строгий поиск не покрывает новые граничные значения окна равенства.
\item \textbf{Вложения.} Классифицировать орбиты примитивных вложений $K_9$, $K_{10}$ и $\Kalt$ в выбранные решётки ранга 48. Каталог из 19 минимальных надсечений одного фиксированного $K_9$ не классифицирует остальные вложения и окружающие орбиты.
\item \textbf{Кластеры вне решённого семейства.} Исследовать сечения ранга 9 без $\sqrt3E_8$, смешанные кластеры с неколлинеарными слоями и оставшиеся диаметры для сечений ранга 10 больших определителей, а также другие ранги, в частности конструкцию размерности 43. Отдельно может исследоваться фактический радиус конкретного окружающего вложения.
\item \textbf{Внешняя экспертиза.} Передать точные формулировки, программы, доказательства полноты и исходные предпосылки специалисту. Внутренний параллельный аудит не является внешним рецензированием.
\item \textbf{Новизна.} Сопоставить дискретную целочисленную теорему с полной доступной литературой и каталогами. Уже установлено, что достигающий пример $K_9$ принадлежит известному совершенному классу; этого различия надо придерживаться в названии и введении будущей статьи.
\end{enumerate}

\newpage
\appendix
\section{Дополнительные равномерные оценки кластеров ранга 10}\label{app:uniform10}

Следующие результаты относятся к \emph{любому определителю} чётной целой решётки $K$ ранга 10 с $m(K)\ge6$, содержащей $J\cong\sqrt3E_8$. Они дополняют теорему \ref{thm:cluster10}, но не доказывают оптимальность всех кластеров во всём этом семействе.

\subsection{Нормировка и геометрия двумерного слоя}

Подрешётка $J$ автоматически примитивна. Пусть $P$ --- проекция $K$ на $(\spanR J)^\perp$, а
\[
T=P^*=K^*\cap(\spanR J)^\perp,\qquad V=\operatorname{covol}(T).
\]
Вычитание ближайшего представителя $J$ из подъёма ненулевого вектора $P$ и равенство $R_J=8/3$ дают
\[
m(P)=m(T^*)\ge c:=10/3.
\]
Поскольку $d(K)=6561d(P)=6561/V^2$, отношение гарантированной центральной плотности к одноточечной упаковке определителя 54675 равно
\begin{equation}\label{eq:uniform10ratio}
R=\frac5{\sqrt3}\,V s\,g(\beta),\qquad
g(\beta)=(1-\beta/6)^{19}.
\end{equation}
Снова $s$ есть число точек кластера, а $\beta<6$ --- квадрат его диаметра.

Пусть $a$ --- длина примитивного вектора $T$ на некоторой решёточной прямой. Примитивный вектор $T^*$, перпендикулярный этой прямой, имеет длину $a/V$: в базисе, начинающемся с этого вектора $T$, высота второго базисного вектора равна $V/a$. Поэтому
\begin{equation}\label{eq:planarspacing}
\frac Va\le\frac1{\sqrt c}.
\end{equation}
В частности, это верно для кратчайшего ненулевого вектора $T$.

\subsection{Все кластеры с одной проекцией}

\begin{proposition}\label{prop:onefiber10}
Если нетривиальный конечный кластер $S\subset K^*$ целиком лежит в одном слое проекции на $J^*$, то $R<1$. Утверждение равномерно по всем определителям $K$ и всем $\beta<6$.
\end{proposition}

\begin{proof}
Слой является сдвигом двумерной решётки $T$. Сначала пусть точки кластера коллинеарны и $a$ --- примитивный шаг вдоль их прямой. Отрезок, их содержащий, имеет длину не меньше $(s-1)a$ и не больше $\sqrt\beta$. Используя \eqref{eq:planarspacing}, получаем
\[
sV\le\frac{s}{s-1}\sqrt{\beta/c}\le2\sqrt{\beta/c},
\qquad
R\le\sqrt{10\beta}\,g(\beta).
\]
Максимум правой части на $(0,6)$ достигается при $\beta=2/13$. Его квадрат равен
\[
\frac{20}{13}\left(\frac{38}{39}\right)^{38}<1.
\]
Равенство для критической точки следует из $6-\beta=38\beta$, а последнее сравнение выполнено точно.

Для неколлинеарного кластера воспользуемся следующей элементарной оценкой площади:
\begin{equation}\label{eq:planaranticode}
sV\le\frac83\beta.
\end{equation}
Докажем её. Пусть $C$ --- выпуклая оболочка кластера и $N=|C\cap T|\ge s$. Добавление всех решёточных точек оболочки не увеличивает диаметр, равный не более $\sqrt\beta$. Обозначим числа внутренних и граничных точек через $I$ и $B$, так что $N=I+B$. Теорема Пика после перехода от $T$ к $\Z^2$ утверждает
\[
\frac{\operatorname{area}(C)}V=I+\frac B2-1
\ge\frac{N-2}{2}.
\]
Следовательно,
\[
NV\le\frac{2N}{N-2}\operatorname{area}(C).
\]
Если $N=3$, оболочка --- треугольник диаметра не больше $\sqrt\beta$, поэтому её площадь не больше $\sqrt3\beta/4$. Получаем $NV\le3\sqrt3\beta/2<8\beta/3$.

Если $N=4$, оболочка --- треугольник или четырёхугольник. Её площадь не больше $\beta/2$: в четырёхугольнике площадь есть половина модуля векторного произведения диагоналей, обе длины которых не больше $\sqrt\beta$. Тогда $NV\le2\beta<8\beta/3$.

Если $N\ge5$, используем плоское изодиаметрическое неравенство: площадь множества диаметра $d$ не превосходит $\pi d^2/4$. Поэтому
\[
NV\le\frac{N\pi}{2(N-2)}\beta
\le\frac{5\pi}{6}\beta<\frac83\beta,
\]
где достаточно $\pi<16/5$. Это доказывает \eqref{eq:planaranticode}. Константа $8/3$ намеренно не оптимизируется.

Подставляя \eqref{eq:planaranticode} в \eqref{eq:uniform10ratio}, получаем
\[
R\le\frac{40}{3\sqrt3}\,\beta g(\beta).
\]
Максимум достигается при $\beta=3/10$ и равен
\[
\frac4{\sqrt3}\left(\frac{19}{20}\right)^{19}<1,
\qquad
\frac{16}{3}\left(\frac{19}{20}\right)^{38}<1
\]
после возведения в квадрат. Это завершает оба случая.
\end{proof}

Теорема Пика и изодиаметрическое неравенство здесь используются как классические точные геометрические теоремы; численная проверка их не заменяет. Одноточечный кластер рассматривается отдельно: $d(K)\ge54675$ уже даёт $R\le1$.

\subsection{Несколько проекций с коллинеарными слоями}

\begin{proposition}\label{prop:collinear10}
Пусть $S\subset K^*$ имеет не меньше двух различных проекций на $J^*$, $\beta<4/3$, и каждый многоточечный слой кластера коллинеарен. Прямые в разных слоях могут иметь различные направления. Тогда
\[
R\le C_{\mathrm{col}}<1,
\]
где точное значение используемой верхней границы есть
\[
C_{\mathrm{col}}=\sqrt{1620\cdot\frac{16}{117}
\left(\frac{304}{351}\right)^{38}}
=0.969365242080\ldots.
\]
\end{proposition}

\begin{proof}
Количество проекций $M$ удовлетворяет $2\le M\le9$, а $\beta\ge2/3$. Если все слои одноточечные, то $s=M\le9$ и $d(K)\ge54675$ дают
\[
R\le9(8/9)^{19}<1.
\]
Эта оценка также меньше заявленной численной границы.

Иначе положим $a=\sqrt{m(T)}$ и обозначим через $t\ge u$ два наибольших числа точек в слоях. Тогда $t\ge2$, $u\ge1$, $s\le t+(M-1)u$. Коллинеарный слой из $s_i$ точек имеет диаметр $d_i\ge(s_i-1)a$.

Усреднение квадратов четырёх поперечных расстояний между концами двух отрезков даёт вертикальный вклад не меньше $(d_i^2+d_j^2)/4$. Угол между прямыми слоёв при усреднении сокращается; поэтому параллельность не нужна. С учётом расстояния между проекциями имеем
\[
\beta\ge\frac23+\frac{d_i^2+d_j^2}{4},\qquad
a^2\big((t-1)^2+(u-1)^2\big)\le4\Delta,\qquad
\Delta=\beta-\frac23.
\]
Для целых $t\ge u\ge1$, $t\ge2$ справедливо
\begin{equation}\label{eq:fiberinteger}
\frac{t+(M-1)u}{\sqrt{(t-1)^2+(u-1)^2}}
\le\max(M+1,\sqrt2M)\le9\sqrt2.
\end{equation}
Действительно, положим $x=t-1\ge y=u-1\ge0$. При $y=0$ отношение равно $(M+x)/x\le M+1$. При $y\ge1$, $r=x/y\ge1$ оно не больше
\[
\frac{2M-1+r}{\sqrt{1+r^2}}.
\]
Его производная имеет знак $1-(2M-1)r<0$ при $M\ge2$, $r\ge1$, поэтому максимум равен $\sqrt2M$ при $r=1$.

Из \eqref{eq:planarspacing}, \eqref{eq:fiberinteger} и \eqref{eq:uniform10ratio} следует
\[
R\le18\sqrt{5\Delta}\,(1-\beta/6)^{19}.
\]
Максимум правой части достигается при
\[
6-\beta=38(\beta-2/3),\qquad
\beta=\frac{94}{117},\qquad \Delta=\frac{16}{117}.
\]
Его квадрат есть $1620(16/117)(304/351)^{38}<1$, что проверено точной рациональной арифметикой.
\end{proof}

\subsection{Границы результата и сертификат}

Предложение \ref{prop:onefiber10} охватывает все однослойные кластеры при любом $\beta<6$ и любом определителе. Предложение \ref{prop:collinear10} добавляет многослойные кластеры с коллинеарными слоями при $\beta<4/3$. Смешанные кластеры с неколлинеарным двумерным слоем, а также общие многослойные кластеры большего диаметра этими доказательствами не исключены. Суммирование однослойной оценки само по себе недостаточно.

Эти два результата используют конечную таблицу классов $E_8/3E_8$, примитивное двумерное двойственное расстояние и элементарную геометрию. Им не нужна глобальная нижняя оценка $e_9(6)$, теорема $\gamma_9=2$ или большой перебор флагов. Точные сравнения, критические точки и константы проверяет \file{clusters/verify_rank10_single_projection.py}; его JSON-сертификат и журнал находятся рядом.

\clearpage
\begin{thebibliography}{10}
\bibitem{audit}
Внутренний отчёт ветки. \emph{Rank-nine determinant bound: audit and full theory}. 10 сентября 2026. Исходный пакет \file{rank9_audit_evidence_20260910.zip}, русская и английская версии; не опубликован.

\bibitem{source}
Ivan Dorofeev, Xiaoming Sun, Chengu Wang.
\emph{Optimal Extensions of Cross-Sections: Sphere Packings in Dimensions 38 to 43}.
Предоставленный неопубликованный файл \file{main (4).pdf}.

\bibitem{dw}
Mathieu Dutour Sikirić, Wessel van Woerden.
\emph{The lattice packing problem in dimension 9 by Voronoi's algorithm}.
arXiv:2508.20719v2, 7 February 2026.
\url{https://arxiv.org/html/2508.20719v2}.

\bibitem{neumaier}
Arnold Neumaier.
\emph{Bounding basis reduction properties}.
Designs, Codes and Cryptography 84 (2017), 237--259.
Theorem 3.4.
\url{https://doi.org/10.1007/s10623-016-0273-9}.

\bibitem{cs}
J. H. Conway, N. J. A. Sloane.
\emph{Complex and Integral Laminated Lattices}.
Transactions of the American Mathematical Society 280 (1983), 463--490.
\url{https://neilsloane.com/doc/Me106.pdf}.

\bibitem{martinet}
Jacques Martinet, Christian Batut. Каталог решёток и таблицы минимума 6.
\url{https://jamartin.perso.math.cnrs.fr/Lattices/index.html};
\url{https://jamartin.perso.math.cnrs.fr/Lattices/Min6.gp}.

\bibitem{nv}
Gabriele Nebe, Boris Venkov.
\emph{The strongly perfect lattices of dimension 10}.
Journal de Théorie des Nombres de Bordeaux 12 (2000), 503--518.
\url{https://www.numdam.org/item/JTNB_2000__12_2_503_0.pdf}.

\bibitem{pp}
W. Plesken, M. Pohst.
\emph{Constructing Integral Lattices with Prescribed Minimum. II}.
Mathematics of Computation 60 (1993), 817--825.
\url{https://doi.org/10.1090/S0025-5718-1993-1176715-1}.

\bibitem{vwthesis}
Wessel P. J. van Woerden.
\emph{Perfect Quadratic Forms: an Upper Bound and Challenges in Enumeration}.
MSc thesis, Leiden, 2018, sections 10.2--10.3.
\url{https://wesselvanwoerden.com/publication/van-2018-perfect/msc-scriptie-wpj-van-woerden.pdf}.

\bibitem{watson}
G. L. Watson. \emph{The number of minimum points of a positive quadratic form}.
Dissertationes Mathematicae 84 (1971), 1--46.
\url{https://eudml.org/doc/268650}.
Историческая ссылка по \cite{dw}; полный оригинальный PDF в текущем запуске не прочитан.
\end{thebibliography}
\end{document}
